\documentclass[11pt]{article}
\usepackage[a4paper,margin=1in]{geometry}
\usepackage{amsmath,amssymb,amsthm}
\usepackage[hidelinks]{hyperref}
\pdftrailerid{}

\newtheorem{lemma}{Lemma}
\newtheorem{proposition}[lemma]{Proposition}
\newcommand{\Q}{\mathbb Q}
\newcommand{\Gal}{\operatorname{Gal}}

\title{CM Descent for One-Place Archimedean Stark Invariants:\\
Major-Revision Working Draft}
\author{Hainan Zhao}
\date{30 July 2026}

\begin{document}
\maketitle

\begin{abstract}
This is a major-revision working draft, not a submission manuscript
and not a source of proved packet claims.  It isolates the precise
cyclic-quartic normalization used in a proposed descent from
one-place real-quadratic invariants to Stark's proved theorem over
imaginary quadratic fields.  The remaining proof obligations are an
explicit Fourier-to-CM-norm bridge and, for the
\(\Q(\sqrt6)\) example, a prime-ideal and finite-valuation analysis of
the auxiliary Euler factors.  Previously computed packets are retained
as candidates until those obligations close.
\end{abstract}

\section{Status and separation of automorphisms}

Let \(k\) be imaginary quadratic and let \(E/k\) be cyclic quartic,
with
\[
 G=\Gal(E/k)=\langle\sigma\rangle\simeq C_4,\qquad
 \tau=\sigma^2.
\]
The involution \(\tau\) fixes \(k\).  It is therefore \emph{not}
global complex conjugation, which acts nontrivially on \(k\).
Whenever \(E\) is stable under global complex conjugation, denote the
latter by \(j\) in the normal closure and put \(E^+=E^{\langle j\rangle}\).
The internal anti-unit relation concerns \(\tau\); the positive CM
norm concerns \(j\).  They must not be conflated.

Fix a faithful ray character with
\(\psi(\sigma)=i\), and let \(S\) contain the complex place and all
conductor primes.  Suppose Stark's proved imaginary-quadratic
rank-one theorem applies and \(|S|\ge3\), so that it supplies a global
unit \(\varepsilon\in\mathcal O_E^\times\), unique modulo
\(\mu(E)\).  Put \(e=|\mu(E)|\) and
\[
 \ell_g=\log|g\varepsilon|_{\rm ord}.
\]

\begin{lemma}[Cyclic-quartic normalization]\label{lem:norm}
Assume in addition that
\[
 \tau(\varepsilon)\equiv\varepsilon^{-1}\pmod{\mu(E)}.
\tag{1}
\]
Then
\[
 \zeta'_S(0,g)=-\frac2e\ell_g
\tag{2}
\]
and, with \(L_S(s,\psi)=\sum_{g\in G}\overline{\psi(g)}
\zeta_S(s,g)\),
\[
 L'_S(0,\psi)=-\frac4e(\ell_1-i\ell_\sigma).
\tag{3}
\]
These coefficients are independent of the numerical value of \(e\).
\end{lemma}

\begin{proof}
Stark's normalized complex absolute value is
\(|z|_w=|z|_{\rm ord}^2\).  His formula
\(\zeta'_S(0,g)=-e^{-1}\log|g\varepsilon|_w\) gives (2).
Condition (1) gives
\(\ell_{\sigma^2}=-\ell_1\) and
\(\ell_{\sigma^3}=-\ell_\sigma\).  Since
\[
 (\overline{\psi(1)},\overline{\psi(\sigma)},
 \overline{\psi(\sigma^2)},\overline{\psi(\sigma^3)})
 =(1,-i,-1,i),
\]
Fourier inversion of (2) gives (3).
\end{proof}

\begin{lemma}[Torsion invariance]
If \(j\) acts as complex conjugation on roots of unity, then the
ordered magnitudes \((|g\varepsilon|)_{g\in G}\) and positive products
\[
 q_g=(g\varepsilon)\,j(g\varepsilon)
\tag{4}
\]
are unchanged when \(\varepsilon\) is multiplied by an element of
\(\mu(E)\).
\end{lemma}

\begin{proof}
Magnitudes are unchanged by roots of unity.  If
\(\xi\in\mu(E)\), then \(j(\xi)=\xi^{-1}\), so the additional factor
in (4) is \(\xi j(\xi)=1\).
\end{proof}

\section{The missing Fourier-to-norm bridge}

Equality of induced representations proves equality of the relevant
Artin \(L\)-functions, but does not by itself identify an individual
one-place invariant.  The packet results require the following
proposition, which is a proof obligation rather than a theorem of this
draft.

\begin{proposition}[Required bridge; not yet proved]\label{prop:bridge}
Let \(\theta\) be a selected one-place ray character over the real
quadratic field \(K\), and suppose
\[
 \operatorname{Ind}_{G_K}^{G_\Q}\theta
 \simeq
 \operatorname{Ind}_{G_k}^{G_\Q}\psi
\tag{5}
\]
as linear representations, with all imprimitive local Euler factors
specified.  For a ray class \(A\), an explicitly determined
group-ring element \(b_A\in\mathbb Z[G]\) and sign
\(\epsilon_A\in\{\pm1\}\) should satisfy
\[
 \log X_A
 =\epsilon_A\log\left|
 N_{E/E^+}\!\left(\varepsilon^{\,b_A}\right)\right|.
\tag{6}
\]
The map \(A\mapsto b_A\), all local factors, and the Artin labels must
be derived by Fourier inversion rather than inferred from matching
polynomials.
\end{proposition}

To prove Proposition~\ref{prop:bridge}, the revision must print:

\begin{enumerate}
\item both Fourier inversions, including the one-place sign class;
\item the exact character correspondence arising from (5);
\item every local Euler factor at primes removed from either side;
\item the resulting integral group-ring coefficient \(b_A\);
\item the exact normal-closure map taking (4) to the asserted
one-place ray field; and
\item the Artin action on every displayed packet value.
\end{enumerate}

\section{The auxiliary-prime gate for
\texorpdfstring{\(\Q(\sqrt6)\)}{Q(sqrt(6))}}

The current computation establishes two rational-prime Euler
polynomials, normalized analytic runs at \(q=3,5\), exact
unit-lattice divisibility, and the cross-route identity
\(q_8^3=q_{12}^2\).  Those facts do not yet eliminate a residual
\(S\)-unit ambiguity.

Before the case can be promoted, the revision must, for each
imaginary quadratic base and each rational \(q\):

\begin{enumerate}
\item factor \(q\mathcal O_k\) and name every prime
\(\mathfrak q\) added to \(S\);
\item state whether \(\mathfrak q\) is split, inert, or ramified and
whether it divides the conductor;
\item compute \(\psi(\mathfrak q)\) and the group-ring Euler factor
\(1-\operatorname{Frob}_{\mathfrak q}^{-1}N\mathfrak q^{-s}\);
\item compute the finite valuations of the enlarged Stark unit at
all primes above \(q\); and
\item invoke a precise uniqueness theorem showing that equality of
the archimedean logarithms and those finite valuations leaves only a
root of unity.
\end{enumerate}

Only after these checks may a statement such as
``the enlarged unit is the square of the natural unit modulo
torsion'' enter a proof.

\section{Candidate corpus, without promotion}

The existing computations for the \(\Q(\sqrt{35})\) packet, the first
\(e=6\) tranche, RQ-000458, and \(\Q(\sqrt6)\) remain useful candidate
data and regression controls.  They are not claimed as theorems in
this draft.  A future submission will require
Proposition~\ref{prop:bridge} and, for \(\Q(\sqrt6)\), the complete
auxiliary-prime gate above.

\end{document}
